\chapter*{Capstone: Cleaning a Real Messy Dataset}
\label{chap:capstone}
\addcontentsline{toc}{chapter}{Capstone: Cleaning a Real Messy Dataset}
\markboth{Capstone: Cleaning a Messy Dataset}{}

Every earlier lab handed you data that was already tidy. This one does not.
\texttt{Unclean Dataset.csv} is a student enrolment file that has been through
several careless hands, and the task is to turn it into something a database
could accept --- then say honestly what the cleaning cost.

\textbf{Source files.} \texttt{Lab Final Project/Lab Final Project.R},
\texttt{Unclean Dataset.csv} (136 data rows), and
\texttt{Cleaned Dataset.csv} (the expected 64-row result).\\
\textbf{Goal.} Repair a file with mixed delimiters, shifted columns, inconsistent
categories, three date formats, and currency text, and report what was
discarded.

\begin{labmode}
Run it from the folder that holds it, so the relative paths resolve:

\begin{lstlisting}[style=dsr]
setwd("path/to/DATA SCIENCE/Lab Final Project")
source("Lab Final Project.R")
\end{lstlisting}

It needs \texttt{tidyverse} and \texttt{lubridate}. The script checks for them
and stops with instructions rather than installing them, because a lab script
that silently writes to a shared machine's package library is bad manners.
\end{labmode}

\section*{Seven defects, and what each one teaches}

\begin{dstable}
\begin{tabularx}{\linewidth}{@{}>{\bfseries\leavevmode\color{DSNavy}}p{0.21\linewidth}X@{}}
\toprule
\rowcolor{DSPaleBlue!92}
Defect in the file & Why it cannot simply be read with \texttt{read.csv()} \\
\midrule
Two delimiters & Some rows use \texttt{|}, others use \texttt{,}. One
\texttt{read\_csv()} call cannot parse both, so the script splits the lines
first and parses each group with the right delimiter. \\
Trailing commas & Pipe rows end in \texttt{\$1200,,,,,,,}, which would create
phantom empty columns. \\
Shifted columns & In some comma rows the age has slipped into the
\texttt{Gender} field, so a value like \texttt{M24} holds two variables at once. \\
Inconsistent categories & Gender appears as \texttt{M}, \texttt{Male},
\texttt{male}, \texttt{F}, \texttt{Female}; the script folds case and keys on
the first letter. \\
Misspelt courses & \texttt{Machine Learnin}, \texttt{Web Develpment}. Repaired
by pattern, not by a lookup table, because the truncations vary. \\
Three date formats & \texttt{ymd}, \texttt{dmy}, and \texttt{mdy} in one column.
\texttt{parse\_date\_time()} is given all three orders. \\
Currency as text & \texttt{\$1,200} is a string. Strip every non-digit, then
convert. \\
\bottomrule
\end{tabularx}
\end{dstable}

\begin{keyidea}
Notice the order the script works in. It repairs \emph{structure} first (split
the delimiters, unshift the columns), then \emph{values} (categories, dates,
currency), and only then drops rows. Reversing that order throws away rows that
were merely mis-parsed rather than genuinely bad.
\end{keyidea}

\section*{The number that matters most}

The script now prints a cleaning summary. On the supplied file it reads:

\begin{lstlisting}[style=dsr]
raw data lines read      : 136
  pipe-delimited         : 76
  comma-delimited        : 60
rows after merge         : 136
rows after clean/dedupe  : 64
dropped                  : 72
Dates that failed to parse: 7
\end{lstlisting}

\begin{pitfall}
\textbf{53\% of the rows did not survive.} A cleaning script that printed only
\texttt{head(clean\_df)} would look like a complete success while discarding
more than half the dataset, and nobody reading the tidy output would know.

Before you present the cleaned file, account for those 72 rows. How many were
true duplicates removed by \texttt{distinct()}? How many failed the
\texttt{!is.na(Student\_ID)} filter, and were those genuinely unusable or just
parsed badly? Seven enrolment dates are still \texttt{NA} even in the output
you keep.

``Rows dropped'' is not a footnote. It is the main finding of a cleaning
exercise, and it belongs in the report next to the row count you kept.
\end{pitfall}

\begin{tryit}
Add one line before the \texttt{distinct()} call to count exact duplicates
separately from rows removed by the \texttt{is.na} filter:

\begin{lstlisting}[style=dsr]
cat("exact duplicates:", nrow(full_df) - nrow(distinct(full_df)), "\n")
\end{lstlisting}

Then inspect the seven unparsed dates and decide whether a fourth format order
would rescue them, or whether they are genuinely unreadable:

\begin{lstlisting}[style=dsr]
full_df$Enrollment_Date[is.na(clean_df$Enrollment_Date)]
\end{lstlisting}

Either answer is acceptable in the report. Not looking is not.
\end{tryit}

\begin{keyidea}
The script also adds a \texttt{System\_Record\_ID} surrogate key, because
\texttt{Student\_ID} repeats in this file. A natural key that is not actually
unique will break the first join you attempt, so a cleaning step that will feed
a database should end by giving every row an identifier it controls.
\end{keyidea}

\begin{labdeliverable}
Submit the cleaned CSV, the script, and a short data-quality note stating: the
original and final row counts, how many rows each cleaning rule removed and
why, which fields you standardised and to what convention, how many dates
remain unparsed, and one sentence on whether the surviving 64 rows are still
representative of the original file.
\end{labdeliverable}
